\section{Discussion and policy implications}\label{sec:discussion}

\paragraph{Who benefits?} The special issue asks who benefits from AI-driven climate solutions. Our
answer for agriculture, the sector in which heat reaches the poor most directly, is that the
benefit is not yet visible in aggregate outcomes. Countries whose climate-AI research capacity
exceeds what their income would predict suffer the same yield losses from hot months as other
countries. This does not mean that AI-based adaptation tools fail where they are deployed. It means
that national research capacity, the most widely available measure of AI capability, has not
translated into economy-wide resilience large enough to detect, and that claims about AI's
contribution to adaptation should be judged against measured damages rather than technical
potential. The finding is consistent with \citet{larosa2024}, who argue that AI helps vulnerable
communities only when it is accessible to them. Research output that remains in universities, or
that is produced mainly by foreign partners, may not reach the extension services, insurers and
farmers through whom resilience would be built.

\paragraph{Under what conditions?} The heterogeneity estimates are imprecise, but their pattern is
informative for policy design. The coefficient on inequality suggests that capability is, if
anything, less protective in more unequal countries, mirroring the finding of \citet{dong2024}
that inequality weakens AI's carbon-abatement benefits. The learned allocation rules, although unstable across specifications, repeatedly
select on two conditions, connectivity and knowledge sovereignty, rather than on income or existing
AI strength. Poorly connected countries in hot climates, where yield damages are large and
agricultural exposure is high, appear in every rule. Where connectivity is higher, the rules favour
countries whose climate-AI research is predominantly domestic. Both conditions point to the same
policy lever: AI capacity yields resilience only if it is embedded in local institutions and if
the people exposed to heat can be reached.

\paragraph{Who bears the costs?} The economy-wide energy and emissions footprint of national AI
research capacity is too small to detect. This is not evidence that AI is environmentally benign.
The electricity and water demands of AI arise in data centres \citep{iea2025,devries2023}, which are
concentrated in a few, mostly high-income locations, and the countries that build research capacity
for adaptation are not those that host the infrastructure. The distributional question on the cost
side is therefore a cross-border one, between the places that bear the environmental footprint of
AI infrastructure and the places that might benefit from AI-enabled adaptation. It cannot be
answered with national research indicators.

\paragraph{Implications for emerging and developing economies.} Three implications follow for
governments and development partners in emerging and developing economies, including Viet Nam and
its ASEAN neighbours. First, the most common allocation heuristic, building on existing strength,
performs worst in our evaluation. Concentrating capacity where climate-AI research is already
strongest yields almost no equity-weighted resilience, because those countries have little
agricultural exposure. Second, programmes should be evaluated against measured damages. The
method in this paper can be applied to programme-level or subnational data, where both damages
and capability are measured more precisely than in cross-country aggregates. Third, the choice of
inequality aversion is a political decision that our framework makes explicit. In our application
it changes the recommended set of countries only at high values of $\epsilon$. At $\epsilon=2$ the
rule includes Viet Nam and four other ASEAN members, because the region combines high exposure of
lower-income rural households to heat with modest domestic climate-AI research. For Viet Nam, where
agriculture remains an important source of income for poorer rural households, the result supports investment in domestically led applied
climate-AI, such as agro-advisory and early-warning services, rather than in research output as
such.

\paragraph{Limitations.} Four limitations qualify these conclusions. First, bibliometric output
measures research capacity, not deployment. A direct measure of AI use in agriculture and disaster
management would sharpen the estimates considerably, and no such measure currently exists across
countries. Second, overlap between countries with and without income-adjusted capability is
limited, and about half of country-years have propensity scores outside $[0.1,0.9]$. This is a
property of the data, not of the estimator, but it limits precision. Third, identification of the
resilience effect rests on selection on observables for shock sensitivity
(Assumption~\ref{as:select}). Unobserved factors that raise both research capacity and agricultural
resilience, such as strong agricultural research institutes, could bias the estimates in either
direction. Fourth, the decile exposure of agricultural income is approximated from employment
shares. Household-survey data on income sources by decile would allow a direct measure.
